% Companion long-record migration for Erdős #269 round 6.
% Destination label: long:earlier-ordinary-results (historical / two-dimensional remainder of original sec:open)
% Not frozen R. Not a blind \input. Labels may collide with the short note;
% assign destination labels as specified in R6LongRecordMigrations.md.
% Extracted from Type B edits/relocation_extracts.tex.

\subsection{A two-dimensional analytic representation}

For the de-duplicated series define
\[
 \mathcal D_{2,3,5}=1+
 \sum_{t\in\{2^n,3^n,5^n:n\ge1\}}
 \frac{1}{
   2^{\lfloor\log_2t\rfloor}
   3^{\lfloor\log_3t\rfloor}
   5^{\lfloor\log_5t\rfloor}}.
\]
Irrationality of $\mathcal D_{2,3,5}$ is the historical de-duplicated
assertion of the 1973 letter, not an open problem of this note.  The problem
below asks only for a function-level representation or a proof that none
exists in a specified class.
The dyadic coding is the joint rotation word
\[
 \delta_{3,a}=\lfloor(a+1)\theta_3\rfloor-\lfloor a\theta_3\rfloor,
 \quad \theta_3=\frac{\log2}{\log3},\qquad
 \delta_{5,a}=\lfloor(a+1)\theta_5\rfloor-\lfloor a\theta_5\rfloor,
 \quad \theta_5=\frac{\log2}{\log5},
\]
with $\beta_a=2\,3^{\delta_{3,a}}5^{\delta_{5,a}}$.

\begin{problem}[function-faithful two-dimensional representation]
Express $\mathcal D_{2,3,5}$ as a nonconstant algebraic combination of values
of a specified two-dimensional Hecke--Mahler, cone-generating or multivariate
Mahler function and verify every hypothesis of a published value theorem; or
give a conditional theorem under an explicit logarithmic nondegeneracy
hypothesis; or prove that the literal series has no representation in the
specified finite-dimensional class.
\end{problem}

Pairwise irrationality of $\theta_3$ and $\theta_5$ is not silently promoted
to the orbit-closure or equidistribution hypothesis a two-dimensional theorem
may need.  The finite-observer formalisation isolates the precise
faithfulness requirement: equality in a finite observer must imply equality
after symbolic realisation, and a genuine finite-dimensional factorisation
forces the realised symbolic span to be finite-dimensional
(\lword{Erdos269/WeightedPhaseCarry.lean}{109}{carry_eq_residueDigit_add_coboundary}{residue--coboundary form},
\lword{Erdos269/WeightedPhaseCarry.lean}{293}{CartierRealisation}{symbolic realisation},
\lword{Erdos269/WeightedPhaseCarry.lean}{334}{finite_realisedSpan_of_factorisation}{checked}).
The same finite formulation keeps the carry residue and residue digit in their
declared intervals
(\lword{Erdos269/WeightedPhaseCarry.lean}{150}{carryResidue_mem_interval}{carry interval},
\lword{Erdos269/WeightedPhaseCarry.lean}{157}{residueDigit_mem_interval}{digit interval}).
No theorem here proves that the literal realised span is infinite, so a
general finite separable decomposition is neither assumed nor declared
excluded.

\subsection{Further structural constraints}

The radix alphabet itself extends without difficulty.  For ordered primes
$p_1<\cdots<p_s$, an interval $(p_1^a,p_1^{a+1})$ contains at most one power
from each other channel, because consecutive $p_i$-powers have ratio
$p_i>p_1$.  Hence its block radix belongs to the $2^{s-1}$-letter alphabet
\[
 \left\{p_1\prod_{i=2}^s p_i^{\varepsilon_i}:
   \varepsilon_i\in\{0,1\}\right\}.
\]
The nontrivial question is quantitative: prove effective recurrence or
discrepancy for the actual four-letter $\{2,6,10,30\}$ word, an asymptotic
formula with an error term for the restricted two-dimensional shell counts
that generate $m_a^{235}$, or the exact finite-separation rank of the literal
three-prime kernel under a specified family of shifts.

The three-channel rigidity and carry-lift extinction theorems already exclude
one proposed argument in four exact steps.  Under channel surjectivity, ordinary
block-nullity is equivalent to the perturbation being a coboundary of a
channel potential
(\lword{Erdos269/ThreeChannelBlockRigidity.lean}{59}{channelBlockNull_iff_channelPotential}{potential
classifier}).  Zero perturbations on genuine $2\to3$ and $2\to5$ transitions
then identify all three potential values and force the perturbation to vanish
at every index.  For an integral lift, that vanishing makes a nonzero initial
lift error grow by the exact product of the successive bases; bases at least
two make its absolute value at least $2^N$, contradicting even a single
index-$N$ bound strictly below $2^N$
(\lword{Erdos269/CarryLiftExtinction.lean}{178}{no_carryLift_of_errorBound_below_twoPow}{one-index
extinction}).  A single index therefore already suffices, and consequently no
uniform bound on the lift error can hold either: under the same hypotheses a
nonzero initial error is incompatible with any bound valid at every index
(\lword{Erdos269/CarryLiftExtinction.lean}{238}{no_bounded_carryLift_of_blockNull_twoAnchors}{uniform
extinction}).

A separate four-state calculation reaches the obstruction earlier: four real
states in $(0,1)$ with unit-accuracy integral lifts and the two anchor
equalities force the first complete $2$-block sum to be $1$, hence that block
cannot be null
(\lword{Erdos269/CarryLiftExtinction.lean}{289}{binaryLift_twoAnchors_force_firstBlock_sum_one}{first-block
sum},
\lword{Erdos269/CarryLiftExtinction.lean}{308}{no_unitAccurateLift_with_twoAnchors_and_firstTwoBlockNull}{four-state
obstruction}).

These lift obstructions assume the two anchor equalities and block-nullity.
The actual orbit and its ordered-power word have already been constructed;
the specified faithful lift and its additional hypotheses have not.  The affine alternative above may produce an
arbitrary integral state, not necessarily zero; its cofinal $1/31$ separation
is neither eventual nor positive-density and gives no unbounded distance.
The four-state calculation supplies no cofinal windows, and
\texttt{rational\_of\_scaledTail\_integer} classifies no denominators.  The
actual carry supplies a weighted block defect instead, so any successful
argument must use that weighted identity or construct a different faithful
lift.  None of these results is an irrationality theorem for \#269.

The bridge is proved.  The unresolved arithmetic is the nonintegrality
of $BX_a$ for every $a\ge1$ and every positive $B$ coprime to $30$.
